\documentclass[11pt]{article}
\usepackage[margin=1in]{geometry}
\usepackage{amsmath,amssymb,amsthm}
\usepackage{booktabs}
\usepackage{parskip}
% xeCJK is loaded only to typeset the verbatim Chinese statement of the
% conjecture in the quotation below.  A macOS CJK font is used; the rest of the
% document uses only standard LaTeX fonts.
\usepackage{xeCJK}
\setCJKmainfont{PingFang SC}

\newtheorem{theorem}{Theorem}
\newtheorem{lemma}[theorem]{Lemma}
\newtheorem{corollary}[theorem]{Corollary}
\newtheorem{proposition}[theorem]{Proposition}
\theoremstyle{definition}
\newtheorem{definition}[theorem]{Definition}
\theoremstyle{remark}
\newtheorem{remark}[theorem]{Remark}

\newcommand{\Fix}{\operatorname{Fix}}
\newcommand{\wt}{\operatorname{wt}}

\title{Disproof of conjecture \texttt{00000003837}}
\author{earthking11}
\date{2026-09-13}

\begin{document}
\maketitle

\section*{Verdict: FALSE}

We disprove conjecture \texttt{00000003837} as stated. The refutation is
unconditional and elementary. Its first clause fails in the smallest possible
case, type $A_1$ with $\lambda = \omega_1$: the crystal $B(\omega_1)$ has two
elements and the star involution has $0$ fixed points, which is even. Its second
clause fails there too, because $0 < 3$; and in type $A_1$ the hypothesis
$\lambda = -w_0\lambda$ holds for \emph{every} weight, while the fixed-point
count over all of type $A_1$ is at most $1$.

\section{The conjecture}

The statement, quoted verbatim from \texttt{conjectures/00000003837.md}:

\begin{quote}
\textbf{English.} Definition: The star involution $*$ denotes the crystal
involution reversing all simple arrow directions, always existing in finite
type. Conjecture: The number of fixed points of $*$ in $B(\lambda)$ is always
odd, and when $\lambda = -w_0\lambda$ the number of fixed points is at least 3.
(star fixed point parity)
\end{quote}

\begin{quote}
\textbf{中文。} 定义：星对合 $*$ 指同时反转所有简单箭头方向所得的晶体对合，有限型上恒存在。猜想：$B(\lambda)$ 中 $*$ 的不动点个数恒为奇数，且当 $\lambda = -w_0\lambda$ 时不动点个数不少于 3。（星不动点奇偶）
\end{quote}

The claim has two clauses.

\begin{enumerate}
\item \textbf{Parity clause.} For every dominant integral weight $\lambda$ in
  finite type, the star involution $*$ on $B(\lambda)$ has an odd number of
  fixed points.
\item \textbf{Three clause.} Whenever $\lambda = -w_0\lambda$, that number of
  fixed points is at least $3$.
\end{enumerate}

Both are false, and both fail in the same minimal example.

\section{The elementary parity lemma}

\begin{lemma}[Parity lemma]\label{lem:parity}
Let $\sigma$ be an involution of a finite set $X$. Then
\[
\#\Fix(\sigma) \;\equiv\; |X| \pmod 2 .
\]
\end{lemma}

\begin{proof}
For every $x \in X \setminus \Fix(\sigma)$ we have $\sigma(x) \ne x$, and
$\sigma^2 = \mathrm{id}$ shows that $\sigma$ maps the two-element set
$\{x, \sigma(x)\}$ to itself and exchanges its elements. These $2$-cycles
partition $X \setminus \Fix(\sigma)$, so $|X| - \#\Fix(\sigma)$ is even. Hence
$\#\Fix(\sigma) \equiv |X| \pmod 2$.
\end{proof}

Since $|B(\lambda)| = \dim V(\lambda)$, the lemma yields the correct general
congruence
\begin{equation}\label{eq:general}
\#\Fix(*) \;\equiv\; |B(\lambda)| \;=\; \dim V(\lambda) \pmod 2 .
\end{equation}
In particular the parity clause can hold only for those $\lambda$ with
$\dim V(\lambda)$ odd; as an unconditional assertion about all $\lambda$ it is
false. Applied to $B(\omega_1)$, whose cardinality is $2$, the lemma already
forces $\#\Fix(*)$ to be even.

\section{The counterexample: type $A_1$, $\lambda = \omega_1$}

\begin{definition}
In type $A_1$ (that is, for $\mathfrak{sl}_2$) let $\lambda = \omega_1$. The
crystal is the two-element set
\[
B(\omega_1) = \{\,u_+, u_-\,\}, \qquad
\wt(u_+) = +1, \quad \wt(u_-) = -1,
\]
with a single arrow
\[
f_1(u_+) = u_-, \qquad f_1(u_-) = 0, \qquad
e_1(u_-) = u_+, \qquad e_1(u_+) = 0 .
\]
The star involution $*$ reverses all arrows.
\end{definition}

\begin{theorem}\label{thm:counterexample}
On $B(\omega_1)$ the arrow-reversing involution is forced to be the swap
\[
*(u_+) = u_-, \qquad *(u_-) = u_+,
\]
and it has no fixed point:
\[
\#\Fix(*) = 0 .
\]
Consequently $\#\Fix(*)$ is even, and $\#\Fix(*) < 3$.
\end{theorem}

\begin{proof}
Reversing the unique arrow means $f_1 \circ * = * \circ e_1$ (with $*$ extended
by $*(0) = 0$). Evaluating at the bottom element $u_-$, where $e_1(u_-) = u_+$
and $f_1$ is defined, gives
\[
*(u_+) \;=\; *(e_1(u_-)) \;=\; f_1(*(u_-)) ,
\]
so $*(u_+)$ lies in the image of $f_1$, namely $\{u_-\}$. Hence $*(u_+) = u_-$;
applying the involution law $*^2 = \mathrm{id}$ gives $*(u_-) = u_+$. Both
elements are moved, so $\#\Fix(*) = 0$. As $0$ is even and $0 < 3$, both clauses
fail.
\end{proof}

\begin{corollary}
The hypothesis of the three clause holds at $\lambda = \omega_1$: in type $A_1$
the longest element $w_0 = s_1$ acts on the weight lattice by $-1$, so $-w_0$
acts as the identity and
\[
-w_0\,\omega_1 = \omega_1 .
\]
Nevertheless $\#\Fix(*) = 0 < 3$. Thus the three clause fails although its
hypothesis is satisfied, and the parity clause fails at the same time.
\end{corollary}

\begin{remark}
Nothing here is a boundary or degeneracy artefact. The weight $\omega_1$ is the
simplest non-trivial dominant weight in the simplest finite type, and
$B(\omega_1)$ is the simplest non-trivial crystal. Any faithful definition of
the star involution must agree with the swap by
Theorem~\ref{thm:uniqueness} below.
\end{remark}

\section{The general family $B(k\omega_1)$ and the table}

For $k \ge 1$ the crystal $B(k\omega_1)$ is the $(k+1)$-element chain with
weights
\[
k,\; k-2,\; \dots,\; -k .
\]
Index its elements by positions $0, 1, \dots, k$ from the top, so that $j$ has
weight $k - 2j$. The involution $j \mapsto k - j$ is the unique arrow-reversing
involution, and it negates weights. Its fixed points solve $j = k - j$, i.e.
$2j = k$: there is exactly one when $k$ is even (the weight-$0$ element) and
none when $k$ is odd. Therefore
\begin{equation}\label{eq:closed}
\#\Fix(*) \;=\;
\begin{cases}
1, & k \text{ even},\\[2pt]
0, & k \text{ odd}.
\end{cases}
\end{equation}
For $k = 1, \dots, 5$:

\begin{center}
\begin{tabular}{r l l c c c c}
\toprule
$k$ & $\lambda$ & weights of $B(k\omega_1)$ & $|B(k\omega_1)|$
    & $\#\Fix(*)$ & parity & $\#\Fix(*) \ge 3$? \\
\midrule
$1$ & $\omega_1$  & $1,\,-1$                & $2$ & $0$ & even & no ($0<3$) \\
$2$ & $2\omega_1$ & $2,\,0,\,-2$            & $3$ & $1$ & odd  & no ($1<3$) \\
$3$ & $3\omega_1$ & $3,\,1,\,-1,\,-3$       & $4$ & $0$ & even & no ($0<3$) \\
$4$ & $4\omega_1$ & $4,\,2,\,0,\,-2,\,-4$   & $5$ & $1$ & odd  & no ($1<3$) \\
$5$ & $5\omega_1$ & $5,\,3,\,1,\,-1,\,-3,\,-5$ & $6$ & $0$ & even & no ($0<3$) \\
\bottomrule
\end{tabular}
\end{center}

In the $k = 2$ row the crystal has $3$ elements (odd cardinality) and the
fixed-point count is $1$, consistent with \eqref{eq:general} and odd as the
parity clause would demand --- but $1 < 3$. So the failure of the three clause
is not a mere parity artefact either.

\section{The clause ``$\#\Fix(*) \ge 3$'' fails for every $\lambda$ in type $A_1$}

By \eqref{eq:closed}, over all $k \ge 1$ the fixed-point count of the star on
$B(k\omega_1)$ is either $0$ or $1$; hence
\[
\max_{k \ge 1} \#\Fix\bigl(* \text{ on } B(k\omega_1)\bigr) \;=\; 1 \;<\; 3 .
\]
Since in type $A_1$ one has $-w_0 = \mathrm{id}$, so that
$\lambda = -w_0\lambda$ for every dominant weight $\lambda$, the three clause
\emph{fails for every $\lambda$ in type $A_1$}, not only for $\lambda =
\omega_1$. The parity clause, meanwhile, holds exactly for the even $k$ (for
which $\dim V(k\omega_1) = k+1$ is odd), in accordance with \eqref{eq:general};
this is a tautology on the restricted set of $\lambda$ with $\dim V(\lambda)$
odd, not the unconditional claim of the conjecture.

\section{Convention robustness}

One might try to rescue the conjecture by adopting a different convention for
$*$ or for weights. This is impossible.

\begin{theorem}\label{thm:involution}
Every involution of a two-element set has $0$ or $2$ fixed points. In
particular every involution $\sigma$ of $B(\omega_1)$ has even
$\#\Fix(\sigma)$, and $\#\Fix(\sigma) < 3$.
\end{theorem}

\begin{proof}
A two-element set is $\{a,b\}$. An involution either fixes both elements
(identity, $2$ fixed points) or swaps them (no fixed points). Both counts, $0$
and $2$, are even and less than $3$.
\end{proof}

\begin{theorem}[Uniqueness of the star]\label{thm:uniqueness}
Let $\sigma : B(\omega_1) \to B(\omega_1)$ be an involution reversing the arrow,
i.e. $f_1 \circ \sigma = \sigma \circ e_1$ (with $\sigma(0)=0$). Then $\sigma$
is the swap of Theorem~\ref{thm:counterexample}. Hence any faithful convention
for $*$ agrees with it, and $\#\Fix(*) = 0$ is convention-independent.
\end{theorem}

\begin{proof}
Since $e_1(u_-) = u_+$ and $\sigma^2 = \mathrm{id}$,
\[
\sigma(u_+) = \sigma(e_1(u_-)) = f_1(\sigma(u_-)),
\]
so $\sigma(u_+) \in \{u_-\}$, i.e. $\sigma(u_+) = u_-$ and then
$\sigma(u_-) = u_+$ by involutivity.
\end{proof}

\begin{proposition}\label{prop:no-weight-preserving}
There is no involution of $B(\omega_1)$ that both reverses the arrow and
preserves weights. The forced star negates weights, so the weight condition
cannot be relaxed to manufacture a fixed point.
\end{proposition}

\begin{proof}
By Theorem~\ref{thm:uniqueness} any arrow-reversing involution is the swap,
which sends the top element $u_+$ of weight $+1$ to $u_-$ of weight $-1 \ne +1$.
\end{proof}

Thus the value $\#\Fix(*) = 0$ cannot be evaded by changing the definition of
$*$: the arrow-reversing requirement pins it down, and no involution of
$B(\omega_1)$ at all has an odd or $\ge 3$ fixed-point count.

\section{Literature context}

The genuine results surrounding the star involution are not parity theorems.
The relevant theorem is Stembridge's and Berenstein--Kirillov's statement that
the number of \emph{self-evacuating tableaux} of a shape equals the number of
\emph{domino tableaux} of that shape (J.~R.~Stembridge, \emph{Duke Math.~J.}
\textbf{82} (1996); A.~Berenstein and A.~Kirillov, \emph{Discrete Math.}
\textbf{225} (2000), and references therein). This is a bijective
equinumerosity statement joining evacuation and domino tableaux; it contains no
assertion that $\#\Fix(*)$ is odd. The correct general congruence in this
circle of ideas is the elementary \eqref{eq:general},
\[
\#\Fix(*) \equiv \dim V(\lambda) \pmod 2 ,
\]
which is a consequence of the $2$-cycle pairing of Lemma~\ref{lem:parity}. The
conjecture appears to be a garbled rendering of these results: its parity clause
conflates ``$\dim V(\lambda)$ is odd'' with ``always odd'', and its three clause
has no counterpart in the literature. We state this only as context; the
refutation above is self-contained and does not rely on it.

\section{Reproducibility}

The checks are carried out by \texttt{reproduce.py}, which uses the Python~3
standard library only:

\begin{quote}\ttfamily
\$ python3 reproduce.py\\
... \textit{(the parity lemma, the } \(B(\omega_1)\) \textit{ computation, the
table for } \(k = 1..10\),\textit{ and the robustness checks)}\\
PASS: all checks verified; conjecture 00000003837 is FALSE.
\end{quote}

It models $B(k\omega_1)$ for $k = 1, \dots, 10$ as an $(k+1)$-element chain,
verifies that $j \mapsto k-j$ is an arrow-reversing involution negating weights,
counts its fixed points, asserts that the count is even for every odd $k$ (in
particular that $k=1$ gives $0$), asserts $0 < 3$, verifies the parity lemma
$\#\Fix \equiv |B(k\omega_1)| \pmod 2$, and exhaustively verifies that every
involution of the two-element set has $0$ or $2$ fixed points and that the
arrow-reversing involution of $B(\omega_1)$ is unique. It prints \texttt{PASS}
and exits $0$ exactly when every check holds.

\section{Formalisation}

The counterexample is formalised in the accompanying Lean~4 project under
\texttt{lean4/}. It builds with \texttt{lake build} and is audited by
\texttt{lake env lean Check.lean}, using core Lean only (no Mathlib) and no
\texttt{sorry}. The crystal $B(\omega_1)$ is modelled as \texttt{Fin 2} with
explicit operators, and the main formal statements are

\begin{quote}\ttfamily
theorem star\_reverses\_f : \(\forall\) j, f1 (star j) = (e1 j).map star\\
theorem wt\_star : \(\forall\) j, wt (star j) = - wt j\\
theorem star\_fixedCount : fixedCount star = 0\\
theorem conjecture\_00000003837\_false :\\
\hspace*{2em}fixedCount star = 0 \(\wedge\) fixedCount star \% 2 = 0\\
\hspace*{2em}\(\wedge\) fixedCount star < 3\\
theorem involution\_fixedCount\_even (\(\sigma\)) (h\(\sigma\) : \(\forall\) x, \(\sigma\) (\(\sigma\) x) = x) :\\
\hspace*{2em}fixedCount \(\sigma\) \% 2 = 0\\
theorem arrow\_reversing\_involution\_eq\_star (\(\sigma\)) (h\(\sigma\)) (hf) : \(\sigma\) = star\\
theorem no\_weight\_preserving\_arrow\_reversal :\\
\hspace*{2em}\(\neg\) \(\exists\) \(\sigma\), (\(\forall\) x, \(\sigma\) (\(\sigma\) x) = x) \(\wedge\)
  (\(\forall\) j, wt (\(\sigma\) j) = wt j)\\
\hspace*{4em}\(\wedge\) (\(\forall\) j, f1 (\(\sigma\) j) = (e1 j).map \(\sigma\))
\end{quote}

\noindent Here \texttt{star} is the swap on \texttt{Fin 2}, \texttt{fixedCount}
counts the fixed points over \texttt{List.finRange 2}, and the theorems
\texttt{involution\_fixedCount\_even} and \texttt{arrow\_reversing\_involution\_eq\_star}
formalise the convention-robustness argument. No \texttt{sorryAx} occurs; the
only axioms reported are \texttt{propext} and \texttt{Quot.sound}. See
\texttt{lean4/README.md} for the full statement table and the faithfulness note.

\end{document}
